% GENERATED FILE. Edit canonical lessons, then run npm run generate:books.
% canonical-source-hash: fnv1a64:9846b3392de2ef0a
% canonical-lessons: FR-C70-dejeuner, FR-C70-diner, FR-C70-petit-dejeuner, FR-C70-riz, FR-C70-poire

\chapter{Lunch, Dinner, Breakfast, Rice, Pear}
\label{ch:fr-lunch-dinner-breakfast-rice-pear}
\hlchaptermodality{70}

\begin{chapteropening}
\textbf{By the end of this chapter:} \emph{I can name lunch, dinner, breakfast, rice, a pear.}

Complete the last lesson of chapter 70: I can name lunch, dinner, breakfast, rice, a pear.
\end{chapteropening}

\section[le déjeuner]{le déjeuner — lunch}

\label{lesson:FR-C70-dejeuner}



\begin{quote}
Before the new one: say the French for honey, then the French for oil.
\end{quote}

\subsection*{You'll want to know: le déjeuner}
\textbf{le déjeuner} — \textquotedblleft{}lunch\textquotedblright{}.

\begin{cousinweb}
\textbf{déjeuner} $\leftarrow$ Vulgar Latin \emph{*disieiunare}, \textquotedblleft{}to break the fast\textquotedblright{} — inherited. (EtymDB 2.1 records the Latin verb.) The French word for \textquotedblleft{}dinner\textquotedblright{} comes from the same verb by another path.
\end{cousinweb}

\begin{grammarlens}[title={What you've built}]
One more word for this chapter.
\end{grammarlens}

\subsection*{Your turn}
\begin{itemize}
\raggedright
  \item \emph{Say it:} \emph{le déjeuner}
  \item \emph{Say it:} \emph{le déjeuner}, once more
  \item \emph{Say it:} say \emph{le déjeuner}
  \item \emph{Recall:} say the French for honey, then the French for oil, then say \emph{le déjeuner} again
\end{itemize}

\begin{tcolorbox}[breakable,colback=teal!4,colframe=teal!35!black,title={Before you move on}]
What does le déjeuner mean? (\textquotedblleft{}Lunch\textquotedblright{}.) Say it once more.
\end{tcolorbox}
\section[le dîner]{le dîner — dinner}

\label{lesson:FR-C70-diner}



\begin{quote}
Before the new one: say the French for oil, then the French for lunch.
\end{quote}

\subsection*{You'll want to know: le dîner}
\textbf{le dîner} — \textquotedblleft{}dinner\textquotedblright{}.

\begin{cousinweb}
\textbf{dîner} $\leftarrow$ Old French \emph{disner} $\leftarrow$ Vulgar Latin \emph{*disiunare}, from \emph{disieiunare}, \textquotedblleft{}to break the fast\textquotedblright{} — inherited. (EtymDB 2.1) It is the same verb as \emph{déjeuner}, worn down differently.
\end{cousinweb}

\begin{grammarlens}[title={What you've built}]
One more word for this chapter.
\end{grammarlens}

\subsection*{Your turn}
\begin{itemize}
\raggedright
  \item \emph{Say it:} \emph{le dîner}
  \item \emph{Say it:} \emph{le dîner}, once more
  \item \emph{Say it:} say \emph{le dîner}
  \item \emph{Recall:} say the French for oil, then the French for lunch, then say \emph{le dîner} again
\end{itemize}

\begin{tcolorbox}[breakable,colback=teal!4,colframe=teal!35!black,title={Before you move on}]
What does le dîner mean? (\textquotedblleft{}Dinner\textquotedblright{}.) Say it once more.
\end{tcolorbox}
\section[le petit-déjeuner]{le petit-déjeuner — breakfast}

\label{lesson:FR-C70-petit-dejeuner}



\begin{quote}
Before the new one: say the French for lunch, then the French for dinner.
\end{quote}

\subsection*{You'll want to know: le petit-déjeuner}
\textbf{le petit-déjeuner} — \textquotedblleft{}breakfast\textquotedblright{}. Literally \textquotedblleft{}little lunch\textquotedblright{}.

\begin{cousinweb}
\textbf{petit-déjeuner} — built from the French word for \textquotedblleft{}little\textquotedblright{} and \emph{déjeuner}: the small first breaking of the fast, as English \emph{breakfast} is.
\end{cousinweb}

\begin{grammarlens}[title={What you've built}]
One more word for this chapter.
\end{grammarlens}

\subsection*{Your turn}
\begin{itemize}
\raggedright
  \item \emph{Say it:} \emph{le petit-déjeuner}
  \item \emph{Say it:} \emph{le petit-déjeuner}, once more
  \item \emph{Say it:} say \emph{le petit-déjeuner}
  \item \emph{Recall:} say the French for lunch, then the French for dinner, then say \emph{le petit-déjeuner} again
\end{itemize}

\begin{tcolorbox}[breakable,colback=teal!4,colframe=teal!35!black,title={Before you move on}]
What does le petit-déjeuner mean? (\textquotedblleft{}Breakfast\textquotedblright{}.) Say it once more.
\end{tcolorbox}
\section[le riz]{le riz — rice}

\label{lesson:FR-C70-riz}



\begin{quote}
Before the new one: say the French for dinner, then the French for breakfast.
\end{quote}

\subsection*{You'll want to know: le riz}
\textbf{le riz} — \textquotedblleft{}rice\textquotedblright{}.

\begin{cousinweb}
\textbf{riz} $\leftarrow$ Old French \emph{ris} $\leftarrow$ Italian \emph{riso} $\leftarrow$ Latin \emph{oryza}, from Greek \emph{oryza} — borrowed from Italian. (EtymDB 2.1)
\end{cousinweb}

\begin{grammarlens}[title={What you've built}]
One more word for this chapter.
\end{grammarlens}

\subsection*{Your turn}
\begin{itemize}
\raggedright
  \item \emph{Say it:} \emph{le riz}
  \item \emph{Say it:} \emph{le riz}, once more
  \item \emph{Say it:} say \emph{le riz}
  \item \emph{Recall:} say the French for dinner, then the French for breakfast, then say \emph{le riz} again
\end{itemize}

\begin{tcolorbox}[breakable,colback=teal!4,colframe=teal!35!black,title={Before you move on}]
What does le riz mean? (\textquotedblleft{}Rice\textquotedblright{}.) Say it once more.
\end{tcolorbox}
\section[la poire]{la poire — a pear}

\label{lesson:FR-C70-poire}



\begin{quote}
Before the new one: say the French for breakfast, then the French for rice.
\end{quote}

\subsection*{You'll want to know: la poire}
\textbf{la poire} — \textquotedblleft{}a pear\textquotedblright{}.

\begin{cousinweb}
\textbf{poire} $\leftarrow$ Old French \emph{peire} $\leftarrow$ Vulgar Latin \emph{*pira}, from Latin \emph{pirum}, \textquotedblleft{}pear\textquotedblright{} — inherited. (EtymDB 2.1)
\end{cousinweb}

\begin{grammarlens}[title={What you've built}]
One more word for this chapter.
\end{grammarlens}

\subsection*{Your turn}
\begin{itemize}
\raggedright
  \item \emph{Say it:} \emph{la poire}
  \item \emph{Say it:} \emph{la poire}, once more
  \item \emph{Say it:} say \emph{la poire}
  \item \emph{Recall:} say the French for breakfast, then the French for rice, then say \emph{la poire} again
\end{itemize}

\begin{tcolorbox}[breakable,colback=teal!4,colframe=teal!35!black,title={Before you move on}]
What does la poire mean? (\textquotedblleft{}A pear\textquotedblright{}.) Say it once more.
\end{tcolorbox}
